\chapter{Review of the trproj contribution}
\label{chap:contrib-trproj}

This chapter reviews the \texttt{trproj} branch as a whole: what it set out to do, how it is built,
how it relates to the thesis and to the \texttt{iota} branch it sits on, and what a maintainer
would still have to accept on faith. It is a prose chapter: it introduces no new definitions or
theorems of its own and refers to the nodes of Chapters~\ref{chap:proj}, \ref{chap:trexpr},
\ref{chap:trenv}, \ref{chap:typechecker}, \ref{chap:kernel} and \ref{chap:inductive} by cross
reference. Every number below is taken from a file under \texttt{.blueprint-work/understand/} named
at the point it is used, or from a \texttt{grep}/\texttt{git} command run directly against the
\texttt{trproj} checkout at \texttt{20ec229}; none is invented.

\section{Goal and scope}

The model expression type \texttt{VExpr} (\ref{ind:vexpr}) has no projection node. Before this
branch, the clause of \texttt{TrExprS} that should relate a kernel projection \texttt{.proj S i e}
to its model translation was routed through
\begin{quote}
\texttt{def TrProj : ∀ (Γ : List VExpr) (structName : Name) (idx : Nat) (e : VExpr), VExpr → Prop
:= sorry}
\end{quote}
(master, \srcloc{Lean4Lean/Verify/Typing/Expr.lean}{68}) — an \emph{opaque} definition, not merely
an unproved theorem: nothing could be constructed in it or refuted about it. The seven structural
lemmas about it (\texttt{weak'}, \texttt{weak'\_inv}, \texttt{defeqDFC}, \texttt{wf},
\texttt{uniq}, \texttt{instN}, \texttt{instL}) were \texttt{sorry} as well, so every declaration of
\texttt{Lean4Lean.Verify} whose statement merely mentions \texttt{TrExprS} or \texttt{TrExpr}
carried \texttt{sorryAx} by construction. \texttt{trproj}'s stated goal (\texttt{trproj-commission.md};
recorded in \texttt{understand/contrib-trproj.md}(a)) was to give \texttt{TrProj} a real
definition and discharge as many of the seven lemmas as the definition allows, so that
environments containing Lean's prelude — \texttt{Prod.fst}, \texttt{HAdd.hAdd} and everything else
whose body is a \texttt{.proj} node — stop being permanently out of reach of \texttt{TrEnv}.

The branch is 24 commits over \texttt{iota..trproj} (2026-08-26 to 2026-09-10), touching 17 files
with $+2273/-170$ lines; the total footprint against \texttt{master} is 50 files and
$+7763/-363$ (\texttt{understand/contrib-trproj.md}, header table). Per-line blame over
\texttt{.blueprint-work/blame/} tags roughly 2645 of those lines \texttt{P} (trproj-only), against
roughly 5000 \texttt{I} (iota) and the remainder \texttt{M} (master) — the branch sits strictly on
top of \texttt{iota}, reusing its registered-$\iota$-rule machinery (\ref{rule:isdefeq-pat},
\ref{structure:ind-wf}) rather than duplicating it.

\section{Design}
\label{sec:contrib-trproj-design}

The design decision, unchanged since the first commit, is that a projection is not a new syntactic
form but an application of the recursor the kernel already generates for the structure:
$P_i = S.\mathrm{rec}\;(uss\,i)\;ps\;(\lambda x.\,F_i[f_j := P_j\,x])\;(\lambda f_0\dots f_{n-1}.\,f_i)$, with the earlier projection \emph{functions} $P_j$ chained into field $i$'s
motive so that a dependent field type is expressible. The full derivation and the source layout are
Chapter~\ref{chap:proj}'s subject; this section only traces how the pieces fit into the branch's
own story.

\texttt{Theory/Proj.lean} (\srcloc{Lean4Lean/Theory/Proj.lean}{1}, 449 lines, entirely
\texttt{trproj}, no \texttt{sorry}) builds the expansion — \ref{def:proj-field-selector} through
\ref{def:proj-fn}, its type (\ref{def:proj-motive-body}, \ref{def:proj-ty}) and the arity pin
\ref{def:proj-binder-arity} — and proves the substitution metatheory each builder needs, once
against \texttt{Subst} (\ref{def:subst}) rather than once per operator (commit \texttt{d7ec0fa}:
$+101/-304$ lines in the file, 405 lines of combined churn, net $-203$, 652 to 449 lines;
\texttt{claims-trproj.md} \#21 corrects the earlier reports that quoted 405 itself as the net
change). \texttt{TrProjCtor}
(\ref{struct:trprojctor}, \srcloc{Lean4Lean/Verify/Typing/Expr.lean}{90}) is the consumer: an
eight-field \texttt{Prop}-valued structure — \texttt{pat}, \texttt{params\_length}, \texttt{ctor},
\texttt{field\_lt}, \texttt{minor\_arity}, \texttt{major\_ty}, \texttt{fn\_ty}, \texttt{eq}, verified
directly against \srcloc{Lean4Lean/Verify/Typing/Expr.lean}{90-131} — that names the constructor,
the per-field level list \texttt{uss}, the parameter list and the field-type telescope explicitly,
where the first design existentially bound them all and could not even state which constructor a
projection resolves against. \texttt{TrProj} (\ref{def:trproj}) hides that data existentially, and
\texttt{TrExprS.proj} is re-threaded onto it unchanged except for two new arguments,
\texttt{env} and \texttt{U} — \texttt{TrProj} is no longer purely syntactic, because
\texttt{major\_ty} and \texttt{fn\_ty} are \texttt{HasType} side conditions carried inside the
translation relation, a design asymmetry with every other clause of \texttt{TrExprS} that the
branch's own closing questions flag for the maintainer.

Of \texttt{TrProj}'s seven inherited lemmas, five are proved: \ref{thm:trproj-weak},
\ref{thm:trproj-defeqdfc}, \ref{thm:trproj-mono} (new, required by \texttt{TrExprS.mono}),
\ref{thm:trproj-wf}, \ref{thm:trproj-instn} and \ref{thm:trproj-instl} — six declarations, since
\texttt{weak'} and \texttt{weakN} share \ref{thm:trproj-weak}. Two remain open:
\ref{thm:trproj-weak-inv}, blocked on the master admission
\texttt{VEnv.IsDefEqU.weakN\_iff} (\ref{thm:isdefequ-weakn-iff}), and \ref{thm:trproj-uniq}, the
load-bearing one: \texttt{TrExprS.uniq} inducts through it, so the branch's own uniqueness layer
inherits \texttt{sorryAx} from a gap it introduced. \ref{thm:tr-env-proj-defeq}
(\srcloc{Lean4Lean/Verify/Environment/Lemmas.lean}{1021}) is the headline theorem: a
continuous, roughly 130-line tactic proof (confirmed \texttt{sorry}-free by
\texttt{grep -n -w sorry} over the file, which returns nothing) that the recursor expansion of
projection $i$ reduces to \texttt{fields[i]} when the major premise is definitionally a saturated
constructor spine. Its hypotheses are exactly what \texttt{inferProj} and \texttt{reduceProjCore}
actually check — non-mutual, single-constructor, non-indexed — rather than the recursor's raw
telescope split, and that split (\ref{thm:tr-env-structure-rec}) is \emph{derived} from the
kernel's structure facts rather than assumed, closing review finding F07 of the branch's own
2026-09-07 review.

The type-checker side (Chapter~\ref{chap:typechecker}) gets two things. \texttt{inferProj.WF} is
split into \ref{thm:vtc-inferproj-struct} (new, scoped to what \texttt{TrProjCtor} covers) and
\ref{thm:vtc-inferproj} (general, and now honestly documented as not provable as stated); both
remain \texttt{sorry}. The split also \emph{fixed} a latent bug in master's statement: master's
\texttt{inferProj.WF} concluded $\exists\,ty',\,c.\mathtt{TrTyping}\;(.proj\;st\;i\;e)\;ty\;e'\;ty'$
— asserting that $e'$, the translation of the \emph{struct}, is the translation of the
\emph{projection} — with \texttt{ty'} auto-bound; \texttt{trproj} restates it as
$\exists\,e''\,ty',\,\dots\,e''\,ty'$. Separately, \texttt{inductiveReduceRecCore}
(\ref{def:kernel-iota-core}, \srcloc{Lean4Lean/Inductive/Reduce.lean}{75}) extracts the pure
$\iota$ step of the \emph{executable} kernel's \texttt{inductiveReduceRec}
(\ref{def:kernel-inductive-reduce-rec}) into its own definition, with the caller reduced to
\texttt{return inductiveReduceRecCore info ls recArgs major}. Diffing master against \texttt{trproj}
directly shows this is byte-for-byte the original body with \texttt{majorIdx} recomputed as
\texttt{rval.getMajorIdx} where the caller had bound \texttt{info.getMajorIdx} — behaviour
preserving, not a semantic change. \ref{thm:vtc-iota-reduce-core} then proves
\texttt{inductiveReduceRecCore.WF} about it: this is $\iota$-consumer content (it exercises the
registered-$\iota$-rule interface \ref{thm:tr-env-iota-rec} of Chapter~\ref{chap:trenv}), landed on
\texttt{trproj} by commit \texttt{6fd8a1d}, and the merge \texttt{20ec229} confirms \texttt{iota}
independently re-landed the same 21 lines.

Two test modules complete the picture, both entirely \texttt{trproj} and both \texttt{sorry}-free:
\texttt{Tests/ProjShape.lean} (185 lines) runs the expansion against the real Lean kernel — 10
structure/field pairs plus two negative controls (a wrong level list, a reflexive structure) —
and \texttt{Tests/ProjInhabit.lean} (597 lines) inhabits \texttt{TrProjCtor} on the model side, on
a plain two-field block and a \texttt{Sigma}-shaped dependent block, through
\texttt{VEnv.addInduct}, with four \texttt{\#guard\_msgs in \#print axioms} pins. \texttt{trproj}
also contributes 15 lines to the two \texttt{iota} specification-test files (11 in
\texttt{Tests/ShapeDecide.lean}, 4 in \texttt{Tests/IotaShape.lean}; verified against the per-line
blame, correcting an earlier note that had counted 4).

\section{Relation to the thesis}

\texttt{Theory/Proj.lean}'s module docstring cites the thesis twice, and both citations are
individually accurate but need to be read for what they are, not for what the commit history
originally claimed. The \emph{shape} of the expansion — a projection defined by applying the
recursor with a chosen motive — is the thesis's $\mathrm{inv}_x : \mathrm{acc}\;x \to \forall y.\, y<x\to\mathrm{acc}\;y$ (typesys.tex, the undecidability-of-conversion section), built from
$\mathrm{rec}_{\mathrm{acc}}$ with a specific, non-constant motive; \texttt{unique.tex:109}
generalises the idiom further than the branch's own notes credit it: ``these
$\mathrm{inv}_i$ projection operators can be defined using the recursor, like we demonstrated for
\texttt{acc}'' (\texttt{understand/claims-trproj.md} \#33). The \emph{dependent typing} of the
motive, $P_i\,e : F_i[f_j := P_j\,e]$, is attributed to the primitive $\Sigma$-projection rule
$\pi_2\,p : \beta[\pi_1\,p/x]$ of \texttt{Wtypes.tex} \S5.1, and the per-field level list
\texttt{uss} to \texttt{axioms.tex} \S2.6.3's remark that the recursor's motive universe is fresh
per use.

The point the task brief and \texttt{thesis-map.md} both insist on is exactly this: \emph{the
thesis has no treatment of Lean's \texttt{Expr.proj} at all}. \texttt{grep -c '\textbackslash proj'}
over every \texttt{.tex} file in \texttt{lean-type-theory} returns zero. Structures, structure
$\eta$ and \texttt{inferProj} are never mentioned. The $\pi_2$ rule the docstring borrows belongs to
\texttt{Wtypes.tex} \S5, the chapter that \emph{replaces} general inductive types by eight
primitives (including $\Sigma$) plus a \texttt{good} predicate and subsingleton eliminators into a
parameterised \texttt{Acc} — a system lean4lean does not formalise at all
(\texttt{thesis-map.md} \S2, table row ``\S5 W-types — nothing''). So the ``dependent motive'' the
branch's commit \texttt{811a52a} titles ``define TrProj by the thesis's dependent motive'' is a
borrowed \emph{idiom} applied to a system the thesis never builds a projection theory for, not an
inherited \emph{result} transported from one. \texttt{Theory/Proj.lean}'s own docstring has since
been softened to call this ``the dependent typing \emph{shape}'' of the thesis's primitive
projections, which is accurate; the commit title that started the confusion was not corrected.
Two further ironies, noted by both \texttt{thesis-map.md} and the branch's own review: $\mathrm{inv}_x$
is introduced in the thesis to prove conversion \emph{undecidable}, and is reused here as the
representation of every structure projection; and \texttt{Wtypes.tex} \S5.1 admits an $\eta$ rule,
$(\pi_1 x, \pi_2 x)\equiv x$, precisely because it has \emph{no} recursors for $\Sigma$ to fall
back on — the exact rule this branch cannot state, because \texttt{IsDefEq} (\ref{def:isdefeq})
has no structure-$\eta$ clause. The three-way disagreement is worth stating plainly: the thesis's
soundness model has $\Sigma$-$\eta$; the real Lean kernel has structure $\eta$
(\texttt{tryEtaStructCore}, \srcloc{Lean4Lean/Verify/TypeChecker/IsDefEq.lean}{227}, still
\texttt{sorry} on both \texttt{master} and \texttt{trproj}); the lean4lean model built by this
branch has neither, and none of the three positions justifies either of the others.

What genuinely \emph{is} a contribution, independent of the borrowed vocabulary, is that
\ref{thm:tr-env-proj-defeq} \emph{derives} the thesis's computation rule $\pi_i(a,b)\equiv a_i$
from the registered $\iota$ rule rather than assuming it — the theorem fires the same
$\iota$-registry machinery \texttt{iota} built for ordinary recursion
(\ref{rule:isdefeq-pat}) and $\beta$-reduces through \ref{def:proj-field-selector}, so the
projection computation is a genuine consequence of the model's existing rules, not a new axiom.

\section{History and churn}

The branch reworked its central design decision twice, each time because the previous version was
diagnosed as wrong before it caused a false theorem, not after. \textbf{M1} (\texttt{f252c3c}
through \texttt{b6a5a38}, 2026-08-26 to 08-28) shipped \texttt{TrProj} with a \emph{free}
existential motive; commit \texttt{7a5e96d} then pinned it to a \emph{constant} motive
$\lambda\_.\,\mathit{fieldTy}$, because on a neutral major premise well-typedness constrains the
motive only on constructor-shaped inputs, so two witnesses with motives agreeing on constructors
but differing on a variable would both satisfy \texttt{TrProj} while not being defeq — making
\texttt{TrProj.uniq} false. Commit \texttt{b6a5a38} then introduced \texttt{TrProjCtor} because the
\emph{previous} statement of \texttt{TrEnv.proj\_defeq} was unprovable: the $\iota$ rule's
constructor name and the spine's constructor name were unrelated. \textbf{M2}
(\texttt{7f62db2}, 2026-09-03) is a documentation-only commit recording a dead end: the
\texttt{rec\_reg} route could not establish that a recursor is a \emph{structure} recursor, because
the $\iota$ key records only the \emph{sum} $\mathit{numMotives}+\mathit{numMinors}+\mathit{numIndices}$.
\textbf{M3} (\texttt{fc9fbfc} through \texttt{f7dabf1}, 2026-09-04) is the decisive rework: commit
\texttt{811a52a} replaced the constant motive with the current chained \emph{dependent} one,
because the constant motive cannot type a dependent field such as \texttt{Sigma.snd}; commit
\texttt{83860e9} turned \texttt{proj\_defeq} into a real theorem, still over the raw telescope
split as a hypothesis. \texttt{f7dabf1} is the commit the project's own 2026-09-07 review was run
against.

\textbf{M4} (\texttt{2a901f4} through \texttt{6fd8a1d}, 2026-09-08) is the review response: eleven
commits, eight of them answering named findings, six of which are worth naming individually:
\texttt{e9fefbe} turns \texttt{TrProjCtor} from a positionally
destructured existential into the current named-field structure (F13) and deletes two zero-use
wrappers (F29); \texttt{e202975} restates \texttt{proj\_defeq} on the kernel's structure facts
instead of the telescope-split hypothesis (F07); \texttt{273cd2c} splits \texttt{inferProj.WF} into
the scoped and general lemmas and documents the general one as not provable as stated (B3/F03);
\texttt{acd6b46} adds the sorry-free inhabitation tests (F19); \texttt{d7ec0fa} re-derives every
builder from \texttt{subst} (F12); \texttt{6fd8a1d} splits \texttt{inductiveReduceRecCore} and
proves \texttt{inductiveReduceRecCore.WF}, the first real consumer of the $\iota$ interface
(answering the $\iota$ half of F09; the \emph{projection} half is still unanswered, \S\ref{sec:contrib-trproj-consumers}).
\textbf{M5} (\texttt{20ec229}, 2026-09-10) is the final merge with \texttt{iota}: trivial, one line
changed, since \texttt{iota} had independently re-landed the M4 kernel-refactor content.

No commit in \texttt{iota..trproj} is a literal \texttt{git revert}; the rework is by replacement.
\texttt{hygiene-review.md} quantifies the cost: non-merge commits on the branch total
$+11439/-3399$ against a net $+7762/-363$ over the whole \texttt{master..trproj} range, so roughly
3000 lines the branch itself wrote were later removed inside the branch — about 27\% of what was
written. The worst offenders by commit count are \texttt{Verify/Environment/Lemmas.lean} (22
commits, $+1734/-678$) and \texttt{Theory/Typing/InductiveParams.lean} (14 commits, over half
rewritten); a whole file, \texttt{Theory/Pattern.lean}, was created over four commits and then
deleted (\texttt{cb920e1}/\texttt{3bbdb22}) and does not exist at the tip. Two pairs of commits
share a patch-id or subject line (\texttt{655dd3f}/\texttt{75ffde9}, \texttt{7a68882}/\texttt{6fd8a1d}),
evidence of re-merging a rebased \texttt{iota} without squashing — the kind of thing an upstream
maintainer would ask to be cleaned up before review. Since the repository's own workflow
squash-merges branches, none of this history reaches upstream if the branch is ever proposed; the
cost is review time inside this project, not permanent record.

\section{Claims audit}

\texttt{understand/claims-trproj.md} independently re-checked every claim
\texttt{contrib-trproj.md} attributes to the branch's own documentation
(\texttt{PR\_trproj.md}, \texttt{TRPROJ\_CONTRIBUTION.md}, \texttt{trproj-commission.md},
\texttt{downstream-asks-commission.md}) against the \texttt{20ec229} checkout, using
\texttt{git grep/diff/show/log/blame} and \texttt{lake env lean} re-elaboration; of 50 claims, 28
came back \textbf{TRUE}, 12 \textbf{PARTIAL}, 8 \textbf{FALSE} and 2 \textbf{UNVERIFIABLE} without a
downstream repo. The headline technical claims all check out exactly as stated: \texttt{TrProj} is
a real definition, \ref{thm:tr-env-proj-defeq} is proved, the sorry census, the axiom cone and the
two test files are as described. What follows is the audit of what does \emph{not} check out
cleanly, restated as a description list; verdicts are the claims file's own.

\begin{description}
\item[\textbf{TRPROJ\_CONTRIBUTION.md — FALSE where cited.}]
It is not a description of this branch: the file is dated 2026-08-28 and documents the
\emph{abandoned} constant-motive design from M1. It quotes a \texttt{TrProj} definition that no
longer exists, advertises \texttt{TrEnv.pats\_iota\_inv} and two \texttt{TrProjCtor} wrappers
deleted in M4, counts 15 \texttt{IOTA-TODO} and 3 \texttt{PROJ-TODO} markers where
\texttt{git grep} at \texttt{20ec229} finds zero, and cites three \texttt{PROJ-TODO} line numbers
one of which now names a \emph{proved} theorem (\ref{thm:tr-env-proj-defeq}). It should not be
handed to a reviewer as a current description.
\item[\textbf{``every lemma about it was sorry on master'' — PARTIAL.}]
\texttt{TrProj.weakN} was already proved on master (a one-line \texttt{simpa} off the sorried
\texttt{weak'}), and \texttt{TrProj.mono} did not exist on master at all. Seven of \texttt{TrProj}'s
eight master-era lemmas were \texttt{sorry}, not eight, and \ref{thm:trproj-mono} is a genuinely new
obligation the branch introduced for itself (required by \texttt{TrExprS.mono}).
\item[\textbf{the constant-motive scope note — FALSE, stale.}]
``The constant motive is correct for non-dependent fields; dependent fields are out of scope'' is
\texttt{TRPROJ\_CONTRIBUTION.md} \S A1 describing the design \texttt{811a52a} replaced eight days
later. The current motive is the chained dependent one, and
\texttt{Tests/ProjShape.lean:120-145} exercises \texttt{Sigma.snd}, \texttt{PSigma} and
\texttt{Subtype}.
\item[\textbf{``the general inferProj.WF keeps its statement'' — FALSE.}]
It does not. Master's version asserted, via an auto-bound implicit, that the translation of the
\emph{struct} is the translation of the \emph{projection}; \texttt{trproj} corrected this to an
existential over the projection's own translation (\S\ref{sec:contrib-trproj-design} above). The
branch fixed a latent statement bug in a master theorem, and \texttt{PR\_trproj.md} itself denies
having touched it.
\item[\textbf{``no new trust'' for TrProj.defeqDFC — PARTIAL.}]
\ref{thm:trproj-defeqdfc}'s computed \texttt{sorryAx} cone includes \texttt{VEnv.WF.patsStrong}
(\ref{thm:wf-patsstrong}), a gap the \emph{same author's} \texttt{iota} branch introduces, not only
master's pre-existing $\Pi$/sort-injectivity gaps (\ref{thm:injectivity-open}). ``No new trust
beyond master'' is the wrong claim; the accurate one is ``no trust beyond what \texttt{iota}
already added.''
\item[\textbf{``DELIVERED, do not re-open'' — PARTIAL, misleading as to reach.}]
\texttt{downstream-asks-commission.md}'s verdict on \texttt{TrEnv.proj\_defeq} is true as to proof
status: the theorem is sorry-free in its own body. It is misleading as to reach: the theorem has
zero consumers inside lean4lean (\S\ref{sec:contrib-trproj-consumers}), and the same document's
further claim that it is ``consumed downstream today'' by \texttt{lean-to-lambdabox} could not be
confirmed — that repo pins \texttt{lean4lean} at exactly \texttt{20ec229} but its current branch
mentions neither \texttt{proj\_defeq} nor an \texttt{of\_trEnv} lemma outside a scratch file.
\item[\textbf{the four upstreaming blockers B1--B4 — FALSE as stated, superseded in substance.}]
\texttt{contrib-trproj.json}'s ``not upstreamable as it stands'' verdict names B4 (does not build
against current master) as unresolved; it is not: \texttt{git merge-base --is-ancestor master
trproj} succeeds. B1/B2's mechanized counterexamples
(\texttt{review-artifacts/counterexamples/cexA.lean}, \texttt{cexB.lean}) no longer elaborate
against the redesigned \texttt{VInductDecl.WF} — the spec changed under them. What genuinely
remains is the residual obligation, not the counterexample: \ref{thm:wf-patsstrong} is still
\texttt{sorry} and gates every strong-system consumer, including this branch's own
\ref{thm:trproj-defeqdfc} and \ref{thm:tr-env-proj-defeq}.
\end{description}

Two findings the claims audit flags as \textbf{underclaimed} are worth carrying into the
assessment. First, because master's \texttt{TrProj} was a \texttt{sorry} \emph{definition}, every
master declaration whose type merely mentioned \texttt{TrExprS}/\texttt{TrExpr} carried
\texttt{sorryAx} by construction; \texttt{claims-trproj.md} \#45 reports 663
\texttt{Lean4Lean.*} declarations whose type mentions these relations, of which 614 are now
\texttt{sorryAx}-free — a change in the trust status of most of \texttt{Verify}, undersold by the
branch's own ``five sorries closed'' framing. Second, the kernel-to-model bridge
(\ref{thm:tr-env-structure-rec}, \texttt{TrEnv.ctor\_arity}, \texttt{TrEnv.pats\_iota\_inv\_shape})
is itself \texttt{sorryAx}-free; all of \ref{thm:tr-env-proj-defeq}'s inherited taint enters through
\texttt{VEnv.IsDefEqU.betaN}, i.e.\ through master's unique-typing layer plus
\ref{thm:wf-patsstrong}, not through anything this branch built.

\section{Proof and code hygiene}

\texttt{hygiene-review.md}'s whole-diff review (\texttt{master..trproj}, $+7762/-363$ over 49
files) and the per-chapter adversarial reviews of \texttt{proj}, \texttt{trexpr}, \texttt{trenv},
\texttt{typechecker}, \texttt{syntax} and \texttt{kernel} agree on the overall picture: this is
substantive, carefully engineered work with a specific, bounded set of defects, not the pattern
hygiene-review.md's own closing section calls ``slop'' (grandiose docstrings over trivial lemmas,
\texttt{simp}-spam, \texttt{decide}/\texttt{omega} papering over unchecked statements). Concretely:

\begin{itemize}
\item \textbf{Statement quality.} \ref{thm:vtc-inferproj}'s own docstring says it is not provable
as stated, yet \texttt{inferType'.WF} consumes it unconditionally
(\texttt{hygiene-review.md} (a), \textbf{HIGH}) — keeping a knowingly-false statement on the
soundness theorem's import path is worse than master's admittedly-also-wrong version, and the
\texttt{typechecker} chapter review independently confirms the same headline weakness. Lower down,
\texttt{VExpr.projTy} (\ref{def:proj-ty}) is documented as ``the kernel's \texttt{inferProj}
result'' but its docstring's own identity to \texttt{projMotiveBody} is never proved, and
\texttt{TrProjCtor} types the projection through \texttt{projMotiveBody} instead — the \texttt{proj}
chapter review confirms this identity has no proof anywhere in the repository.
\item \textbf{Proof brittleness.} \ref{thm:tr-env-proj-defeq}, at 131 lines the branch's longest
proof, and \texttt{inductiveReduceRecCore.WF} (\ref{thm:vtc-iota-reduce-core}), at 124, both lean
heavily on \texttt{omega}/\texttt{decide}: the branch's \texttt{omega} density is $2.0\%$ of code
lines against master's $0.4\%$, and \texttt{decide} $1.1\%$ against $0.2\%$
(\texttt{hygiene-review.md} baseline table), concentrated in exactly these two proofs. Three of
\ref{thm:trproj-weak}, \texttt{instN} and \texttt{instL} each rebuild the same eight-field
\texttt{TrProjCtor} record by hand rather than through a shared lemma — the ``prove once against
\texttt{subst}'' factoring stops at the \texttt{Theory/Proj.lean} builder level and is not carried
up into \texttt{Verify/Typing/Lemmas.lean} (\texttt{trexpr} chapter review, confirmed).
\item \textbf{Documentation volume.} Contributed \texttt{Theory/} code runs $25.3\%$ comment
lines against master's $1.7\%$ in the same directory, with a 56-line module header on the 449-line
\texttt{Proj.lean} and a 27-line docstring on a 5-line definition
(\texttt{Theory/Typing/InductiveParams.lean:351}, iota-side) — \texttt{hygiene-review.md} calls
this ``the single loudest AI-authored-looking signal in the diff,'' while noting the content itself
is substantive (derivations, thesis section numbers, explicit statements of what is not modelled)
rather than restatement.
\item \textbf{Unused declarations.} \texttt{unused-contrib.md}'s original, hand-verified count is
\textbf{31} zero-reference declarations out of 875 contributed (19 iota, 12 trproj), a $3.5\%$
rate, lower than the $11.2\%$ consumer-less rate the same resolver finds in upstream master's own
core. A later mechanical regeneration of the same script instead reports 78 out of 889 ($8.8\%$;
$10.3\%$ iota, $6.2\%$ trproj); the report flags this itself as an artefact, not a new finding —
the regenerated \texttt{decls.tsv} mislabels 28 user-written \texttt{instance}s as \texttt{def},
and the rerun adds 19 further entries never individually re-verified by hand — so 31 is the figure
this chapter relies on. On the trproj side specifically, the confirmed dead declarations are the
two witness theorems \texttt{trProj0}/\texttt{trProjDep1} in \texttt{Tests/ProjInhabit.lean} (built
and checked, but read by nothing), two never-projected test-fixture fields (\texttt{V3.h},
\texttt{Refl.next}), \texttt{projMotiveBody\_zero} and \texttt{instFields\_nil} (\texttt{@[simp]},
never invoked explicitly), and the two open, unreferenced \texttt{sorry}s
\texttt{VEnv.IsDefEq.crDefEq\_of\_induct} and \ref{thm:vtc-inferproj-struct}. (\texttt{VExpr.projTy}
is not on the dead list itself; it is reached only from two \texttt{example}s in
\texttt{Tests/ProjShape.lean}.) The resolver's headline finding is about wiring, not volume:
\ref{thm:tr-env-proj-defeq} itself, the branch's headline theorem, has no consumer beyond an
axiom-profile test (\S\ref{sec:contrib-trproj-consumers}).
\item \textbf{Orphaned master code.} \texttt{trproj}'s generalisation of \texttt{lift'\_inst\_hi}
into \texttt{lift'\_instN\_hi} (\ref{fam:proj-fn-subst} territory in Chapter~\ref{chap:proj}) left
a six-declaration master cluster in \texttt{Theory/VExpr.lean}
(\texttt{lift\_r\_one}, \texttt{Subst.lift\_r\_comm}, \texttt{Subst.trunc},
\texttt{Subst.Depth.\{one,id\}}, \texttt{Subst.Depth}) with no live consumer, without removing it
— \texttt{unused-contrib.md} \S5 calls this ``the one place where the branches made the host repo
measurably worse rather than merely not-better.'' The \texttt{syntax} chapter review independently
confirms only \texttt{consN\_cons} of the four-lemma block it moved is used.
\item \textbf{Test discipline (positive).} Both test modules re-elaborate clean and validate the
design on two independent axes — \texttt{ProjShape.lean} against the real kernel's own
\texttt{addDeclCore}/\texttt{whnf}/\texttt{isDefEq}, \texttt{ProjInhabit.lean} against the model's
\texttt{VEnv.addInduct} — which \texttt{hygiene-review.md} calls ``the strongest evidence against a
`slop' reading of the branch.'' The \texttt{proj} chapter review qualifies this: the two negative
controls are weaker than they look (a \texttt{try}/\texttt{catch} wrapper that passes vacuously on
any internal failure, and a ``some clause fails'' check that does not pin \emph{which} clause), and
neither hand-built environment in \texttt{ProjInhabit.lean} is ever shown well-formed, so the
witnesses show the premises of \ref{struct:trprojctor} are jointly satisfiable, not that they hold
in a well-formed environment.
\item \textbf{Formatting and naming.} Zero added lines exceed 100 columns, zero carry trailing
whitespace, and contributed declarations are \emph{shorter} on average than master's in the same
directories (5.6--7.5 code lines per declaration against master's 8.5--12.0). Naming follows
upstream convention throughout, and the branch's preference for named-field
\texttt{structure ... : Prop} (\texttt{TrProjCtor}, \texttt{AddInduct}, \texttt{TrIndType}) over
master's existential-chain style is a real, if stylistically visible, improvement.
\end{itemize}

\section{Sorry and axiom impact}

Directly verified against the \texttt{20ec229} checkout: \texttt{Verify/Typing/Lemmas.lean} carries
exactly two \texttt{sorry}s, at lines 747 (\ref{thm:trproj-weak-inv}) and 995
(\ref{thm:trproj-uniq}); \texttt{Verify/TypeChecker/InferType.lean} carries exactly two, at lines
398 and 410 (\ref{thm:vtc-inferproj-struct}, \ref{thm:vtc-inferproj}); and
\texttt{Verify/Environment/Lemmas.lean} carries none at all (\texttt{grep -n -w sorry} on the file
is empty), so \ref{thm:tr-env-proj-defeq} really is a complete proof. Against master's projection-side
tally of one \texttt{sorry} definition plus seven \texttt{sorry} lemmas plus
\texttt{inferProj.WF}'s own \texttt{sorry} (nine sites total), the branch closes six (the
definition and five of the seven lemmas) and adds one new one
(\ref{thm:vtc-inferproj-struct}), leaving four open: \ref{thm:trproj-weak-inv},
\ref{thm:trproj-uniq}, \ref{thm:vtc-inferproj-struct} and \ref{thm:vtc-inferproj}
(\texttt{claims-trproj.md} \#7). \texttt{understand/contrib-trproj.md}'s sorry census for the whole
non-\texttt{Experimental} tree, applying the convention of discounting comment and docstring hits,
is master 23, \texttt{iota} 21, \texttt{trproj} 16.

\texttt{git diff master..trproj | grep -E '\textasciicircum+\textbackslash s*(axiom|opaque)'} finds
no new axiom or opaque declaration, and \texttt{native\_decide} does not occur anywhere in the
diff; both facts were re-checked directly against the branch and hold. What the branch does add,
inherited from \texttt{iota}, is a new load-bearing \texttt{sorry}: \ref{thm:wf-patsstrong}
(\srcloc{Lean4Lean/Theory/Typing/EnvLemmas.lean}{334}). Because \texttt{VEnv.WF.orderedStrong} is
built from it via a \texttt{CoeOut} instance, \ref{thm:trproj-defeqdfc} and
\ref{thm:tr-env-proj-defeq} both carry \texttt{sorryAx} for this reason, not for any gap in their
own proofs — the axiom cone of \ref{thm:tr-env-proj-defeq} traces, via
\texttt{VEnv.IsDefEqU.betaN}, to exactly \ref{thm:injectivity-open}'s two members
($\Pi$-injectivity and sort-injectivity, both pre-existing master admissions) plus
\ref{thm:wf-patsstrong}. \ref{thm:trproj-uniq}'s open \texttt{sorry} is the more direct debt: it
propagates through \texttt{TrExprS.uniq} into \texttt{TrExprS.defeqDFC} and every builder of
\texttt{TrExpr} (Chapter~\ref{chap:trexpr}'s \texttt{family:trexpr-builders}), and separately
reaches \ref{thm:vtc-iota-reduce-core}, so the branch's one $\iota$-side proof also carries the
projection side's open gap. Read against master's baseline, the qualitative change is larger than
the sorry count suggests: because \texttt{TrProj} was a \texttt{sorry} \emph{definition} on master,
\texttt{TrExprS} itself was \texttt{sorryAx}-tainted for every caller; \texttt{\#print axioms
Lean4Lean.TrExprS} now returns \texttt{[propext, Classical.choice, Quot.sound]}, and
\texttt{TrProj}/\texttt{TrProjCtor} return \texttt{[propext]} alone (\texttt{claims-trproj.md}
\#45).

\section{Consumers: is proj\_defeq used anywhere?}
\label{sec:contrib-trproj-consumers}

No. \texttt{grep -rn proj\_defeq --include="*.lean" Lean4Lean}, run directly against the checkout,
returns six lines: the theorem's own signature
(\srcloc{Lean4Lean/Verify/Environment/Lemmas.lean}{1021}; its doc-comment above never repeats
the name), two docstring
mentions in \texttt{Verify/Typing/Expr.lean} (lines 96 and 115, both prose explaining
\texttt{TrProjCtor}'s fields, not code that calls the theorem), and three lines in
\texttt{Tests/ProjInhabit.lean} (582, 587, 595) that are a \texttt{\#guard\_msgs in \#print axioms}
assertion about its axiom footprint — an axiom-profile check, not a use of the theorem's
conclusion. There is no in-library consumer.

The theorem's three intended consumers are all still \texttt{sorry}:
\texttt{reduceProjCore.WF} (\ref{thm:vtc-reduceprojcore},
\srcloc{Lean4Lean/Verify/TypeChecker/Reduce.lean}{143}, \texttt{sorry} at line 145 — note that
\texttt{reduceProj.WF} (\ref{thm:vtc-reduceproj}) one declaration below it \emph{is} proved, by
delegating to \texttt{reduceProjCore.WF}, a distinction an earlier internal note got backwards),
\ref{thm:vtc-inferproj-struct} and \ref{thm:vtc-inferproj}. Its premise \texttt{TrProjCtor} is only
ever inhabited by hand, on the two fixtures of \texttt{Tests/ProjInhabit.lean}; the general
construction that would inhabit it from an arbitrary kernel-accepted projection is exactly
\ref{thm:vtc-inferproj-struct}, still open. Nor is \ref{thm:tr-env-proj-defeq} reached from outside
lean4lean: \texttt{lean-to-lambdabox} pins this checkout at exactly \texttt{20ec229}
(\texttt{lake-manifest.json}) but a search of its current branch for \texttt{of\_trEnv} or
\texttt{TrEnv.proj\_defeq} returns nothing outside one scratch file
(\texttt{claims-trproj.md} \#25). The kernel-to-model bridge lemmas the proof itself depends on —
\ref{thm:tr-env-structure-rec}, \texttt{TrEnv.ctor\_arity}, \texttt{TrEnv.pats\_iota\_inv\_shape} —
have, symmetrically, no consumer \emph{other than} \ref{thm:tr-env-proj-defeq} itself. And the
\texttt{AddInduct} hypothesis every one of these theorems is stated about is itself never
discharged: \texttt{addDecl.WF}'s \texttt{inductDecl} case (\ref{thm:add-decl-wf}) is still
\texttt{sorry} at \srcloc{Lean4Lean/Verify/Environment.lean}{208}, and no declaration anywhere in
the repository constructs an \texttt{AddInduct} witness from
\texttt{Environment.addInductive} — confirmed by grepping every \texttt{AddInduct} occurrence in
the tree, all of which eliminate it, none of which introduce it. So \ref{thm:tr-env-proj-defeq} is,
today, a real theorem about a hypothesis nobody has yet built.

This is the one place where the $\iota$ half of the branch's own work fares better:
\ref{thm:vtc-iota-reduce-core} \emph{does} have a real consumer relationship to
\ref{thm:tr-env-iota-rec}, the registered-$\iota$-rule interface \texttt{iota} built, discharging
half of review finding F09 (``prove at least one internal consumer''). The \emph{projection} half
of that same finding — a real consumer of \ref{thm:tr-env-proj-defeq} — remains unanswered.

\section{Upstream situation}

Nothing in this branch has ever been proposed upstream. \texttt{upstream-comparison.md} confirms
that no upstream commit, PR or issue in the relevant window touches inductives, $\iota$ or
projections beyond a documentation note (\texttt{3adf6da6}, ``docs: justify the projection
divergence by the generated recursor'') and an unrelated checker performance change
(\texttt{62441418}, \texttt{lazyDeltaProjReduction}); \texttt{git log} on upstream's default branch
in the same window shows no commit by the maintainer that touches
\texttt{Verify/Typing/Expr.lean} or \texttt{Verify/Typing/Lemmas.lean}'s \texttt{TrProj} lemmas.
The files this branch adds outright — \texttt{Theory/Proj.lean} (449 lines),
\texttt{Tests/ProjInhabit.lean} (597), \texttt{Tests/ProjShape.lean} (185), and the
\texttt{TrProj}/\texttt{TrProjCtor} rewrite in \texttt{Verify/Typing/Expr.lean} — exist in no
upstream pull request. \texttt{trproj} is 52 commits ahead of \texttt{master} and 24 ahead of
\texttt{iota}, and \texttt{git merge-base --is-ancestor master trproj} succeeds, so a textual merge
would go through; that is a low bar, since \texttt{master}'s own state at the relevant commit still
has \texttt{AddInduct} as a constructorless \texttt{inductive} and \texttt{TrProj := sorry}, so
almost anything merges cleanly against a placeholder.

\texttt{upstream-comparison.md}'s size argument applies here as much as to \texttt{iota}: at
$+7762/-363$ over 49 files, \texttt{trproj} (together with the \texttt{iota} it depends on, at
$+5628/-331$ over 45 files) sits in the size band the maintainer has previously pushed back on for
unrelated PRs, and it changes three existing interfaces at once
(\texttt{VEnv}'s arity, \texttt{TrProj}'s arity, \texttt{addDecl.WF}'s case split). Set against
that: this is, by the hole-closing count, the more valuable of the two branches per line changed —
it retires the \texttt{TrProj} \texttt{sorry} definition plus five of its seven lemma
\texttt{sorry}s — and its closing questions to a maintainer (quoted in
\texttt{understand/contrib-trproj.md} (e)) are specific and well posed: whether \texttt{TrProj}'s
typing side conditions are acceptable inside a translation relation, whether a real
\texttt{VExpr} projection node with its own $\iota$/$\eta$ rules would be preferred instead, and
whether \texttt{TrProjCtor}'s scope restriction (non-mutual, single-constructor, non-recursive,
non-indexed) is the right cut. \texttt{trproj} cannot be upstreamed before \texttt{iota} is, since
it depends on \texttt{iota}'s registered-$\iota$-rule interface throughout; and \texttt{iota}'s own
open blockers — \ref{thm:wf-patsstrong} chief among them — are therefore also blockers for this
branch, inherited rather than introduced.

\section{Assessment}

\textbf{What is solid.} The design decision is right for this model and cheap in the sense that
matters most: it adds no rule to \texttt{IsDefEq} (\ref{def:isdefeq}), so none of the model's
existing metatheory has to be redone, and a projection computes because the structure's own
$\iota$ rule computes. \ref{thm:tr-env-proj-defeq} is a genuine, complete, non-circular proof of
the thesis-adjacent computation rule, derived rather than postulated. The two design corrections
(free motive, then constant motive, then chained dependent motive) each closed a real defect
identified \emph{before} it produced a false theorem, which is the right way to make mistakes in a
verification project. The test discipline — real-kernel validation in
\texttt{Tests/ProjShape.lean}, sorry-free model-side inhabitation in \texttt{Tests/ProjInhabit.lean}
— is better than most of the surrounding repository, upstream included. No new axiom, no
\texttt{native\_decide}, no line over 100 columns, contributed proofs shorter on average than
master's.

\textbf{What is open.} \ref{thm:trproj-uniq} is the single highest-value remaining item: it is the
one \texttt{sorry} this branch owns that gates \texttt{TrExprS.uniq} and, through it, a large
fraction of \texttt{Verify}. \ref{thm:vtc-inferproj-struct} and \ref{thm:vtc-inferproj} are the
bridge from a real kernel projection to a \texttt{TrProjCtor} witness, and neither exists; until one
does, \ref{thm:tr-env-proj-defeq} is evidence about a satisfiable specification, not a theorem
about anything the checker does. And \ref{thm:add-decl-wf}'s \texttt{inductDecl} case being
\texttt{sorry} means every result in this chapter and in Chapter~\ref{chap:trenv} that mentions
\texttt{AddInduct} is conditional on a hypothesis nothing yet constructs — the branch, like
\texttt{iota} beneath it, has built the specification and its consequences without yet building the
thing that would make either fire on a real inductive declaration.

\textbf{What is weak.} \texttt{TRPROJ\_CONTRIBUTION.md} is stale to the point of being actively
misleading and should not circulate. The thesis attribution in the earliest commit titles overclaims
what \texttt{Theory/Proj.lean}'s current, more careful docstring already corrects. Documentation
volume in the contributed \texttt{Theory/} files is out of proportion with upstream's own norms,
though not with its substance. And \texttt{lift'\_inst\_hi}'s generalisation left a small cluster
of master code newly dead without removing it — a real, if minor, regression in tidiness that an
upstream reviewer would flag on sight.

\textbf{Priority order for closing the open ends.} (1) \ref{thm:trproj-uniq} — it is this branch's
own load-bearing gap, its residual (unique typing of the projection function, type-former
injectivity, functionality of the $\iota$ registry) is precisely named in the theorem's own
docstring, and closing it removes \texttt{sorryAx} from the uniqueness layer of the entire
translation. (2) \ref{thm:add-decl-wf}'s \texttt{inductDecl} case — without a producer for
\texttt{AddInduct}, neither this chapter's theorems nor Chapter~\ref{chap:trenv}'s apply to any
concrete environment, and this is shared, prerequisite work for \texttt{iota} as much as for
\texttt{trproj}. (3) \ref{thm:vtc-inferproj-struct} — inhabiting it in general, even under the
scope restriction \texttt{TrProjCtor} already imposes, is what would let a real kernel projection
ever reach \ref{thm:tr-env-proj-defeq}. (4) \ref{thm:wf-patsstrong} — not owned by this branch, but
the single admission that taints \ref{thm:trproj-defeqdfc} and \ref{thm:tr-env-proj-defeq} alike,
and closing it benefits every chapter in this development, not only this one. (5) Housekeeping:
retire \texttt{TRPROJ\_CONTRIBUTION.md} or mark it superseded, restore a consumer for the six
declarations \texttt{lift'\_inst\_hi}'s move orphaned, and rebase away the duplicated commits before
proposing anything upstream.
